ГРИГОРЬЕВ: (ударяя Семенова по морде):
Вот вам и зима настала. Пора печи топить.
Как по"=вашему?

СЕМЕНОВ:
По"=моему, если отнестись серьезно к вашему замечанию, то,
пожалуй, действительно, пора затопить печку.

ГРИГОРЬЕВ: (ударяя Семенова по морде):
А как по"=вашему, зима в этом году будет холодная или теплая?

СЕМЕНОВ:
Пожалуй, судя по тому, что лето было дождливое, то зима всегда холодная.

ГРИГОРЬЕВ: (ударяя Семенова по морде):
А вот мне никогда не бывает холодно.

СЕМЕНОВ:
Это совершенно правильно, что вы говорите, что вам не бывает холодно. У вас такая натура.

ГРИГОРЬЕВ: (ударяя Семенова по морде):
Я не зябну.

СЕМЕНОВ:
Ох!

ГРИГОРЬЕВ: (ударяя Семенова по морде):
Что ох?

СЕМЕНОВ: (держась за щеку):
Ох! Лицо болит!

ГРИГОРЬЕВ:
Почему болит?

(И с этими словами опять Семенова по морде).

СЕМЕНОВ: (падая со стула):
Ох! Сам не знаю!

ГРИГОРЬЕВ: (ударяя Семенова ногой по морде):
У меня ничего не болит.

СЕМЕНОВ:
Я тебя, сукин сын, отучу драться! (пробует встать).

ГРИГОРЬЕВ: (ударяя Семенова по морде):
Тоже, учитель нашелся!

СЕМЕНОВ: (валится на спину):
Сволочь, паршивая!

ГРИГОРЬЕВ:
Ну ты, выбирай выражения полегче!

СЕМЕНОВ: (силясь подняться):
Я, брат, долго терпел. Но хватит.

ГРИГОРЬЕВ: (ударяя Семенова каблуком по морде):
Говори, говори! Послушаем.

СЕМЕНОВ: (валится на спину):
Ох!

\begin{flushright}
	$<\dots>$
\end{flushright}